\subsection{Basis of Cryptography}

Consider the archetypical case in 
cryptography.
Alice wants to send a message to Bob,
and wants the contents to remain 
secret from Vincent. Alice encrypts 
her message with a secret key $s$.

An encryption scheme (Gen, Enc, Dec) 
is a triple of algorithms where
\begin{itemize}
    \item Key generation Gen($1^n$)
    outputs a secret key $k$ of length 
    $n$ (also called the security parameter).
    \item Encryption algorithm 
    $\text{Enc}_k(m) = c$ outputs ciphertext 
    $c$. 
    \item Decryption algorithm $\text{Dec}_k(c) = m$
    outputs plaintext $m$.  
\end{itemize}
An encryption algorithm is defined over 
a message space $M$ and a key space $K$. 
Requirements:
\begin{itemize}
    \item Correctness: $\text{Dec}_k(\text{Enc}_k(m)) = m$.
    \item Security: definition with respect to adversaries. 
\end{itemize}

The secret key $k$ should be hidden from the 
adversary. The encryption and decryption 
algorithms themselves should be public; this 
is Kerckhoffs's principle. This is because 
the secrey keys are easy to hide and replace 
in case of a breach, not so with algorithms. 
It is also good to have many eyes on an 
algorithm to expose security flaws. 

The key space must be large enough 
such that a brute force attack is 
infeasible. Consider a Caesar cipher;
the key space is only 26 long and a
message so encrypted is quickly bruteforced. 

Shannon's Treatment: messages come 
from some distribution, let $M$ be 
a random variable for sampling the 
messages from the message space $M$. 
Distribution $M$ is known to the 
adversary. This captures a priori 
information about the messages. 
The ciphertext $c = \text{Enc}_k(m)$
depends on $m$ chosen according to $M$. 
$k$ chosen randomly according to $\text{Gen}$. 
Enc may also use some randomness. Dec is 
deterministic. These induce a distribution $C$ 
over the ciphertexts $c$. 
The adversary only observes $c$. 
Knowledge about $m$ before observing 
the output of $C$ is captured by $M$. 
Knowledge about $m$ after observing the 
output of $C$ is captured by $M|C$. 
An encryption scheme (Gen, Enc, Dec) is 
\emph{Shannon secure} if for every probability 
distribution for $M$, every message $m \in M$,
and every ciphertext $c \in C$: 
$Pr[M=m|C=c] = Pr[M=m]$. 

An alternative defition of secrecy is 
\emph{Perfect Secrecy}. An encryption
scheme is perfectly secure for every 
pair of messages $m_1$ and $m_2 \in M$
and every ciphertext $c\in C$ if 
$Pr[\text{Enc}_k(m_1)=c]=Pr[\text{Enc}_k(m_2)=c]$.

Shannon had fingers in many cryptographic pies. 
Shannon's Theorem says that if 
(Gen, Enc, Dec) is a perfectly secret 
encryption scheme with message space $M$ 
and key space $K$, then $|K| \geq |M|$. 
This means that for an encryption scheme 
to be perfectly secret, the keys must 
be as long as the message. The challenge of 
modern cryptography is to make something 
just as good as perfect for practical purposes. 